\section{Kendall calibration with ties}
\label{app:null-asymptotics}

This appendix gives the algebra behind the Kendall calibration described in
Section~\ref{subsec:standard-tests}.  There are two distinct statements.  For
uniform masses, the variance conditional on the two observed marginal tie
patterns is an exact permutation identity.  For unequal masses, the software
inserts fractional pseudo-counts into that identity.  The latter construction
has the correct first-order null limit under diffuse exogenous weighting, but
is not an exact finite-sample permutation variance.

\subsection{Exact conditional variance for uniform masses}

For $n\geq3$, let the non-singleton $X$-tie groups have sizes
$t_1,t_2,\ldots$ and let the
non-singleton $Y$-tie groups have sizes $u_1,u_2,\ldots$.  Write

\begin{align*}
  S
  &= \sum_{1\leq i<j\leq n}
     \operatorname{sgn}(X_i-X_j)\operatorname{sgn}(Y_i-Y_j).
\end{align*}

Under $H_0$, condition on both marginal multisets and assign the observed
$Y$ values uniformly over the fixed $X$ positions.  The conditional mean of
$S$ is zero and its exact variance is \citep{kendall1945,kendall1947}

\begin{align}
  V_{S,0}
  ={}&
  \frac{1}{18}
  \left\{
    n(n-1)(2n+5)
    -\sum_g t_g(t_g-1)(2t_g+5)
    -\sum_h u_h(u_h-1)(2u_h+5)
  \right\}
  \notag\\
  &+
  \frac{
    \left\{\sum_g t_g(t_g-1)\right\}
    \left\{\sum_h u_h(u_h-1)\right\}
  }{2n(n-1)}
  \notag\\
  &+
  \frac{
    \left\{\sum_g t_g(t_g-1)(t_g-2)\right\}
    \left\{\sum_h u_h(u_h-1)(u_h-2)\right\}
  }{9n(n-1)(n-2)}.
  \label{eq:kendall-uniform-tie-variance}
\end{align}

The two tau-b denominator factors are fixed under these permutations.  In the
notation of Section~\ref{sec:standard-measures}, they are

\begin{align*}
  Q_X
  &= \binom{n}{2}-\sum_g\binom{t_g}{2},
  &
  Q_Y
  &= \binom{n}{2}-\sum_h\binom{u_h}{2}.
\end{align*}

Consequently,

\begin{align*}
  \widehat\tau_b
  &= \frac{S}{\sqrt{Q_XQ_Y}},
  &
  \widehat\tau_b
     \sqrt{\frac{Q_XQ_Y}{V_{S,0}}}
  &= \frac{S}{\sqrt{V_{S,0}}}.
\end{align*}

The last statistic has conditional mean zero and conditional variance one
exactly.  Its standard normal calibration is nevertheless an asymptotic
approximation; equation~\eqref{eq:kendall-uniform-tie-variance} does not make
the finite permutation distribution Gaussian.

It is useful to rewrite the same identity in terms of elementary symmetric
sums.  For a tie group $G$ of unit masses,

\begin{align*}
  e_2((1)_{i\in G})
  &= \binom{|G|}{2},
  &
  e_3((1)_{i\in G})
  &= \binom{|G|}{3},\\
  \left\{|G|^2-|G|\right\}(2|G|+5)
  &= |G|(|G|-1)(2|G|+5).
\end{align*}

These are respectively the pair, triplet, and variance tie terms that the
unequal-weight construction replaces by pseudo-count analogues.

\subsection{Fractional pseudo-count adjustment for unequal masses}

Put $N=n_{\mathrm{eff}}$ and rescale the normalized case masses as

\begin{align*}
  q_i
  &= Np_i.
\end{align*}

This scale is characterized by

\begin{align*}
  \sum_iq_i
  &=N,
  &
  \sum_iq_i^2
  &=N,
  &
  e_2(q_1,\ldots,q_n)
  &=\frac{N(N-1)}{2}.
\end{align*}

For a marginal tie group $G$, define its pseudo-count power sums and its pair
and triplet masses by

\begin{align*}
  q(G)
  &=\sum_{i\in G}q_i,
  &
  q_r(G)
  &=\sum_{i\in G}q_i^r,\\
  E_2(G)
  &=e_2((q_i)_{i\in G})
    =\frac{q(G)^2-q_2(G)}{2},\\
  E_3(G)
  &=e_3((q_i)_{i\in G})
    =\frac{q(G)^3-3q(G)q_2(G)+2q_3(G)}{6}.
\end{align*}

For the collections $\mathcal G_X$ and $\mathcal G_Y$ of maximal marginal
tie groups, set

\begin{align*}
  T_{X,2}
  &=\sum_{G\in\mathcal G_X}E_2(G),
  &
  T_{Y,2}
  &=\sum_{G\in\mathcal G_Y}E_2(G),\\
  T_{X,3}
  &=\sum_{G\in\mathcal G_X}E_3(G),
  &
  T_{Y,3}
  &=\sum_{G\in\mathcal G_Y}E_3(G),\\
  T_{X,v}
  &=\sum_{G\in\mathcal G_X}
    \{q(G)^2-q_2(G)\}\{2q(G)+5\},\\
  T_{Y,v}
  &=\sum_{G\in\mathcal G_Y}
    \{q(G)^2-q_2(G)\}\{2q(G)+5\}.
\end{align*}

Also write

\begin{align*}
  E_2
  &=e_2(q_1,\ldots,q_n),
  &
  E_3
  &=e_3(q_1,\ldots,q_n).
\end{align*}

The pseudo-count version of
equation~\eqref{eq:kendall-uniform-tie-variance} is

\begin{align}
  V_{S,q}^{\mathrm{app}}
  ={}&
  \frac{2E_2(2N+5)-T_{X,v}-T_{Y,v}}{18}
  +\frac{T_{X,2}T_{Y,2}}{E_2}
  +\frac{2T_{X,3}T_{Y,3}}{3E_3}.
  \label{eq:kendall-pseudocount-variance}
\end{align}

The concordance numerator on this scale and the weighted tau-b coefficient
are

\begin{align*}
  S_q
  &=\sum_{i<j}q_iq_j
    \operatorname{sgn}(X_i-X_j)\operatorname{sgn}(Y_i-Y_j),\\
  \widehat\tau_w
  &=\frac{S_q}{
    \sqrt{(E_2-T_{X,2})(E_2-T_{Y,2})}}.
\end{align*}

Thus the implemented adjustment factor gives

\begin{align}
  Z_K^{\mathrm{app}}
  &=\widehat\tau_w
    \sqrt{
      \frac{(E_2-T_{X,2})(E_2-T_{Y,2})}
           {V_{S,q}^{\mathrm{app}}}
    }
   =\frac{S_q}{\sqrt{V_{S,q}^{\mathrm{app}}}}.
  \label{eq:kendall-pseudocount-statistic}
\end{align}

If the original masses are uniform, then $N=n$, every $q_i=1$,
$E_2=\binom n2$, and $E_3=\binom n3$.  The group terms above reduce to the
integer terms in equation~\eqref{eq:kendall-uniform-tie-variance}, so
$V_{S,q}^{\mathrm{app}}=V_{S,0}$ exactly.

For unequal masses, the equality with a conditional permutation variance no
longer holds.  The $q_i$ need not be integers, literal replication would
introduce within-case ties absent from $S_q$, and permuting responses over
unequal case masses changes the weighted mass of each response tie group and
hence the tau-b denominator.  Equation~\eqref{eq:kendall-pseudocount-variance}
is therefore a scale-invariant finite-sample approximation.  An exact
conditional test with unequal masses must permute the responses and recompute
the complete weighted statistic, including its ranks and denominator, on
every permutation.

Without observed ties, equation~\eqref{eq:kendall-pseudocount-statistic}
reduces to

\begin{align*}
  Z_K^{\mathrm{app}}
  &=\widehat\tau_w
    \sqrt{\frac{9N(N-1)}{2(2N+5)}}
   =\frac{3}{2}\sqrt N\,\widehat\tau_w+o_p(1),
\end{align*}

which agrees with Proposition~\ref{prop:standard-null-limits} to first order.
The software formula requires $E_3>0$, and all displayed studentizations
require positive tau-b denominator factors.

\subsection{First-order null variance with population ties}

The pseudo-count formula is also compatible with a rigorous asymptotic
calculation when ties arise from atoms.  Let $X'$ be an independent copy of
$X$ and define

\begin{align*}
  a_X
  &=\mathbb P(X\neq X'),
  &
  r_X(x)
  &=\mathbb E\{\operatorname{sgn}(x-X')\}
    =F_X(x^-)+F_X(x)-1,\\
  m_X
  &=\mathbb E\{r_X(X)^2\}.
\end{align*}

Define $a_Y$, $r_Y$, and $m_Y$ analogously.

\begin{proposition}[Weighted tau-b under independence with ties]
\label{prop:kendall-population-ties}
Suppose that the observations are independent and identically distributed,
the weights are exogenous and satisfy
equation~\eqref{eq:weight-diffuseness}, $X$ and $Y$ are independent, and
$a_Xa_Y>0$.  Then

\begin{align*}
  \sqrt{n_{\mathrm{eff}}}\,\widehat\tau_w
  &\mathrel{\Longrightarrow}
  N\!\left(0,\frac{4m_Xm_Y}{a_Xa_Y}\right),\\
  \frac{\sqrt{n_{\mathrm{eff}}}\,\widehat\tau_w}
       {2\sqrt{m_Xm_Y/(a_Xa_Y)}}
  &\mathrel{\Longrightarrow}N(0,1).
\end{align*}
\end{proposition}

\begin{proof}
Let

\begin{align*}
  h((x,y),(x',y'))
  &=\operatorname{sgn}(x-x')\operatorname{sgn}(y-y').
\end{align*}

Under independence, its expectation is zero and its first projection is

\begin{align*}
  \mathbb E\{h((x,y),(X',Y'))\}
  &=r_X(x)r_Y(y).
\end{align*}

The weighted Hoeffding decomposition of the concordance numerator, followed
by division by its two tau-b denominator factors, gives

\begin{align*}
  \widehat\tau_w
  &=\frac{1}{\sqrt{a_Xa_Y}}
    \sum_i
    \frac{p_i(1-p_i)}{e_2(p_1,\ldots,p_n)}
    r_X(X_i)r_Y(Y_i)
    +o_p(\sqrt{s_{2,n}})\\
  &=\frac{2}{\sqrt{a_Xa_Y}}
    \sum_i p_i r_X(X_i)r_Y(Y_i)
    +o_p(\sqrt{s_{2,n}}).
\end{align*}

The denominator fluctuations do not enter at first order because the null
concordance numerator is zero.  Moreover, independence gives

\begin{align*}
  \mathbb E\{r_X(X)r_Y(Y)\}
  &=0,\\
  \operatorname{Var}\{r_X(X)r_Y(Y)\}
  &=m_Xm_Y.
\end{align*}

The weighted triangular-array central limit theorem now proves both claims.
\end{proof}

If the atom masses of $X$ are $(\pi_{X,g})_g$, then

\begin{align}
  a_X
  &=1-\sum_g\pi_{X,g}^2,
  &
  m_X
  &=\frac{1-\sum_g\pi_{X,g}^3}{3},
  \label{eq:kendall-atom-functionals}
\end{align}

and similarly for $Y$.  The second identity follows by introducing two
independent copies $X'$ and $X''$:

\begin{align*}
  m_X
  &=\mathbb E\{
    \operatorname{sgn}(X-X')\operatorname{sgn}(X-X'')
  \}.
\end{align*}

If the three values are distinct, the sign product is $1$ for either extreme
and $-1$ for the middle value, so its average over the three labels is $1/3$.
If exactly two values are equal, the product is $1$ for the singleton and
zero for either tied label, again giving average $1/3$.  The product is zero
when all three values are equal.  The probability of the latter event is
$\sum_g\pi_{X,g}^3$, proving
equation~\eqref{eq:kendall-atom-functionals}.

Finally, the leading terms of the pseudo-count summaries explain why the
finite-sample approximation has the same population-tie limit.  Under the
conditions of Proposition~\ref{prop:kendall-population-ties}, diffuseness and
exogeneity imply

\begin{align*}
  \frac{T_{X,2}}{N^2}
  &\mathrel{\longrightarrow_p}
    \frac12\sum_g\pi_{X,g}^2,
  &
  \frac{T_{X,3}}{N^3}
  &\mathrel{\longrightarrow_p}
    \frac16\sum_g\pi_{X,g}^3,\\
  \frac{T_{X,v}}{N^3}
  &\mathrel{\longrightarrow_p}
    2\sum_g\pi_{X,g}^3,
  &
  \frac{E_2}{N^2}
  &\longrightarrow\frac12,
  &
  \frac{E_3}{N^3}
  &\longrightarrow\frac16,
\end{align*}

with analogous limits for $Y$.  Substitution into
equation~\eqref{eq:kendall-pseudocount-variance} yields

\begin{align*}
  \frac{V_{S,q}^{\mathrm{app}}}{N^3}
  &\mathrel{\longrightarrow_p}
  \frac{
    \left(1-\sum_g\pi_{X,g}^3\right)
    \left(1-\sum_h\pi_{Y,h}^3\right)
  }{9},\\
  \frac{E_2-T_{X,2}}{N^2}
  &\mathrel{\longrightarrow_p}\frac{a_X}{2},
  &
  \frac{E_2-T_{Y,2}}{N^2}
  &\mathrel{\longrightarrow_p}\frac{a_Y}{2}.
\end{align*}

Consequently, the factor in
equation~\eqref{eq:kendall-pseudocount-statistic} is asymptotically the same
as the standardization in
Proposition~\ref{prop:kendall-population-ties}.  This agreement is a
first-order result and does not turn the unequal-weight construction into an
exact finite-sample permutation identity.
